% Notas pendientes asociadas a las referencias bibliográficas.
% Se incluyen tras \printbibliography, sin encabezado de capítulo.

\begin{pendiente}[Por completar --- referencias]
\begin{itemize}
    \item Agregar la URL exacta del sitio web de BeeTracer y la fecha de consulta efectiva del
          equipo (entrada \texttt{bilix\_sf} en \texttt{referencias.bib}).
    \item Agregar la fecha exacta de la charla técnica y el nombre del expositor (entrada
          \texttt{charlabeetracer2026}).
\end{itemize}
\end{pendiente}

\begin{pendiente}[Por completar --- recomendado: fuentes del dominio logístico]
Para respaldar el Capítulo~\ref{cap:problema} (definición del problema) conviene citar al menos una
fuente formal del dominio. Se sugiere:

\begin{itemize}
    \item Servicio Nacional de Aduanas de Chile: normativa sobre el manifiesto electrónico y el
          formato XML de transmisión.
    \item Servicio Nacional de Aduanas de Chile: regulación de los almacenes de depósito aduanero y
          terminales extraportuarios.
    \item Empresa Portuaria Valparaíso: estadísticas de transferencia de carga que respalden la
          afirmación sobre la saturación del puerto.
\end{itemize}
\end{pendiente}
